\begin{table*}[t]
\centering
\small
\caption{Synthesis of Eight Core Architectural Decision Dimensions in RouteMate AI.}
\label{tab:tradeoff_synthesis}
\begin{tabularx}{\textwidth}{l p{3.2cm} p{3.2cm} p{3.2cm} p{3.8cm}}
\toprule
\textbf{Dimension} & \textbf{Tension} & \textbf{Low-Cost Baseline} & \textbf{High-Fidelity Regime} & \textbf{Recommended Policy} \\
\midrule
\textbf{Ranking Quality vs Heuristic Simplicity} & Complex supervised models vs transparent rule-based linear indices. & Rule-based weighted index (Exp 001, 004): 0.824 NDCG@3, zero training, instant auditability. & Supervised logistic regression (Exp 005): 0.797 NDCG@3, learns feature weights, requires labeled training data. & Deploy transparent rule-based scoring for candidate retrieval, reserving ML for user-acceptance personalization. \\
\textbf{Road-Network Accuracy vs Euclidean Approximation} & Rapid straight-line geometric approximation vs turn-constrained street graph topology. & Euclidean distance: ~45 us/pair, 0 graph memory, ignores one-way restrictions and turn constraints. & Network-aware Dijkstra: ~974 us/pair (21.5x slower), respects one-way streets, turn delays, and network topology. & Never use raw Euclidean distance for final dispatch; always evaluate true street graph shortest paths. \\
\textbf{Routing Latency vs Exact Graph Search} & Sub-millisecond query requirements vs multi-stop Dijkstra path search overhead. & Straight-line / approximate geometry: ~45 us/call, enables 20,000 evaluations/second. & Full Dijkstra routing on multi-stop pooled routes: 4-15 ms per vehicle permutation (Exp 011). & Enforce early-stopping branch pruning and precomputed contraction-hierarchy distance tables for road networks. \\
\textbf{Pruning Recall vs Search Speedup} & Aggressive heuristic candidate filtering vs provably admissible bounding. & Heuristic circuity pruning (Exp 009): 97.8\% pruned, but 85.7\% false negatives (discards valid matches). & Admissible lower-bound pruning (Exp 009): 81.8\% pruned, 5.49x speedup, exactly 0 false negatives, 100\% recall. & Deploy two-tier admissible Euclidean lower-bound pruning as a mandatory pre-filter before road routing. \\
\textbf{Assignment Optimality vs Solver Runtime} & Global branch-and-bound optimization vs greedy priority queue heuristics. & Greedy insertion: < 1 ms runtime across all tested scales, provable 10-17\% optimality gap. & Exact branch-and-bound: guarantees 0\% gap, but times out or exceeds 1.1s on cohorts >= 20x10. & Use greedy insertion for real-time online dispatch, utilizing exact solvers only for offline benchmark auditing. \\
\textbf{Vehicle Pooling Capacity vs Passenger Detour Burden} & High vehicle passenger capacity vs individual rider travel-time inflation. & Single-occupancy C=1 (Exp 011): 0 rider detour, but only 40.0\% matching fulfillment. & Multi-rider pooling C=2..4 (Exp 011): 68-72\% matching fulfillment (+70\%), but 1.22x rider duration and 2.0x detour. & Standardize on vehicle capacity C=2 for private-car pooling, enforcing a 3.5 km hard detour cap. \\
\textbf{Congestion Delay \& Recourse vs Static Plan Fragility} & Pre-planned static dispatch vs dynamic reactive route adaptation. & Static unadapted execution (Exp 012): 0 online compute overhead, but 40-50\% feasibility failure under closures. & Online recourse (Exp 012): 3.42 ms latency, evaluates alternate sequences, restores feasibility to 100\%. & Equip all active pooled vehicles with online incident recourse routers operating with sub-5ms local replanning. \\
\textbf{Rolling-Horizon Dispatch vs Batch Quantization Delay} & Fixed-interval batching (60s) vs continuous event-driven replanning. & Static batching (Exp 013): fixed 60s windows, no active-trip insertions, Fleet VKT = 127.8 km, p50 wait = 169s. & Event-driven rolling dispatch (Exp 013): immediate trigger on arrivals/stops/cancellations, Fleet VKT = 50.5 km (-60.5\%), p50 wait = 103s (-39.1\%). & Standardize platform dispatch on event-driven rolling horizon with active-trip waypoint insertions. \\
\bottomrule
\end{tabularx}
\end{table*}
